% Options for packages loaded elsewhere
% Options for packages loaded elsewhere
\PassOptionsToPackage{unicode}{hyperref}
\PassOptionsToPackage{hyphens}{url}
\PassOptionsToPackage{dvipsnames,svgnames,x11names}{xcolor}
%
\documentclass[
  letterpaper,
  DIV=11,
  numbers=noendperiod]{scrartcl}
\usepackage{xcolor}
\usepackage{amsmath,amssymb}
\setcounter{secnumdepth}{-\maxdimen} % remove section numbering
\usepackage{iftex}
\ifPDFTeX
  \usepackage[T1]{fontenc}
  \usepackage[utf8]{inputenc}
  \usepackage{textcomp} % provide euro and other symbols
\else % if luatex or xetex
  \usepackage{unicode-math} % this also loads fontspec
  \defaultfontfeatures{Scale=MatchLowercase}
  \defaultfontfeatures[\rmfamily]{Ligatures=TeX,Scale=1}
\fi
\usepackage{lmodern}
\ifPDFTeX\else
  % xetex/luatex font selection
\fi
% Use upquote if available, for straight quotes in verbatim environments
\IfFileExists{upquote.sty}{\usepackage{upquote}}{}
\IfFileExists{microtype.sty}{% use microtype if available
  \usepackage[]{microtype}
  \UseMicrotypeSet[protrusion]{basicmath} % disable protrusion for tt fonts
}{}
\makeatletter
\@ifundefined{KOMAClassName}{% if non-KOMA class
  \IfFileExists{parskip.sty}{%
    \usepackage{parskip}
  }{% else
    \setlength{\parindent}{0pt}
    \setlength{\parskip}{6pt plus 2pt minus 1pt}}
}{% if KOMA class
  \KOMAoptions{parskip=half}}
\makeatother
% Make \paragraph and \subparagraph free-standing
\makeatletter
\ifx\paragraph\undefined\else
  \let\oldparagraph\paragraph
  \renewcommand{\paragraph}{
    \@ifstar
      \xxxParagraphStar
      \xxxParagraphNoStar
  }
  \newcommand{\xxxParagraphStar}[1]{\oldparagraph*{#1}\mbox{}}
  \newcommand{\xxxParagraphNoStar}[1]{\oldparagraph{#1}\mbox{}}
\fi
\ifx\subparagraph\undefined\else
  \let\oldsubparagraph\subparagraph
  \renewcommand{\subparagraph}{
    \@ifstar
      \xxxSubParagraphStar
      \xxxSubParagraphNoStar
  }
  \newcommand{\xxxSubParagraphStar}[1]{\oldsubparagraph*{#1}\mbox{}}
  \newcommand{\xxxSubParagraphNoStar}[1]{\oldsubparagraph{#1}\mbox{}}
\fi
\makeatother


\usepackage{longtable,booktabs,array}
\newcounter{none} % for unnumbered tables
\usepackage{calc} % for calculating minipage widths
% Correct order of tables after \paragraph or \subparagraph
\usepackage{etoolbox}
\makeatletter
\patchcmd\longtable{\par}{\if@noskipsec\mbox{}\fi\par}{}{}
\makeatother
% Allow footnotes in longtable head/foot
\IfFileExists{footnotehyper.sty}{\usepackage{footnotehyper}}{\usepackage{footnote}}
\makesavenoteenv{longtable}
\usepackage{graphicx}
\makeatletter
\newsavebox\pandoc@box
\newcommand*\pandocbounded[1]{% scales image to fit in text height/width
  \sbox\pandoc@box{#1}%
  \Gscale@div\@tempa{\textheight}{\dimexpr\ht\pandoc@box+\dp\pandoc@box\relax}%
  \Gscale@div\@tempb{\linewidth}{\wd\pandoc@box}%
  \ifdim\@tempb\p@<\@tempa\p@\let\@tempa\@tempb\fi% select the smaller of both
  \ifdim\@tempa\p@<\p@\scalebox{\@tempa}{\usebox\pandoc@box}%
  \else\usebox{\pandoc@box}%
  \fi%
}
% Set default figure placement to htbp
\def\fps@figure{htbp}
\makeatother


% definitions for citeproc citations
\NewDocumentCommand\citeproctext{}{}
\NewDocumentCommand\citeproc{mm}{%
  \begingroup\def\citeproctext{#2}\cite{#1}\endgroup}
\makeatletter
 % allow citations to break across lines
 \let\@cite@ofmt\@firstofone
 % avoid brackets around text for \cite:
 \def\@biblabel#1{}
 \def\@cite#1#2{{#1\if@tempswa , #2\fi}}
\makeatother
\newlength{\cslhangindent}
\setlength{\cslhangindent}{1.5em}
\newlength{\csllabelwidth}
\setlength{\csllabelwidth}{3em}
\newenvironment{CSLReferences}[2] % #1 hanging-indent, #2 entry-spacing
 {\begin{list}{}{%
  \setlength{\itemindent}{0pt}
  \setlength{\leftmargin}{0pt}
  \setlength{\parsep}{0pt}
  % turn on hanging indent if param 1 is 1
  \ifodd #1
   \setlength{\leftmargin}{\cslhangindent}
   \setlength{\itemindent}{-1\cslhangindent}
  \fi
  % set entry spacing
  \setlength{\itemsep}{#2\baselineskip}}}
 {\end{list}}
\usepackage{calc}
\newcommand{\CSLBlock}[1]{\hfill\break\parbox[t]{\linewidth}{\strut\ignorespaces#1\strut}}
\newcommand{\CSLLeftMargin}[1]{\parbox[t]{\csllabelwidth}{\strut#1\strut}}
\newcommand{\CSLRightInline}[1]{\parbox[t]{\linewidth - \csllabelwidth}{\strut#1\strut}}
\newcommand{\CSLIndent}[1]{\hspace{\cslhangindent}#1}



\setlength{\emergencystretch}{3em} % prevent overfull lines

\providecommand{\tightlist}{%
  \setlength{\itemsep}{0pt}\setlength{\parskip}{0pt}}



 


% Define chapter counter if document class doesn't have one (e.g. article)
\ifdefined\thechapter\else
  \newcounter{chapter}
  \renewcommand*{\thechapter}{\arabic{chapter}}
\fi

\usepackage{setspace}
\onehalfspacing

\usepackage[pagewise]{lineno}
\linenumbers

\usepackage{orcidlink}

\usepackage{xcolor}
\usepackage{tikz}
\usepackage{afterpage}
\usepackage{eso-pic}
\definecolor{thesiscream}{RGB}{255, 252, 240}
\definecolor{thesissand}{RGB}{232, 212, 162}
\pagecolor{thesiscream}
\KOMAoption{captions}{tableheading}
\makeatletter
\@ifpackageloaded{caption}{}{\usepackage{caption}}
\AtBeginDocument{%
\ifdefined\contentsname
  \renewcommand*\contentsname{Table of contents}
\else
  \newcommand\contentsname{Table of contents}
\fi
\ifdefined\listfigurename
  \renewcommand*\listfigurename{List of Figures}
\else
  \newcommand\listfigurename{List of Figures}
\fi
\ifdefined\listtablename
  \renewcommand*\listtablename{List of Tables}
\else
  \newcommand\listtablename{List of Tables}
\fi
\ifdefined\figurename
  \renewcommand*\figurename{Figure}
\else
  \newcommand\figurename{Figure}
\fi
\ifdefined\tablename
  \renewcommand*\tablename{Table}
\else
  \newcommand\tablename{Table}
\fi
}
\@ifpackageloaded{float}{}{\usepackage{float}}
\floatstyle{ruled}
\@ifundefined{c@chapter}{\newfloat{codelisting}{h}{lop}}{\newfloat{codelisting}{h}{lop}[chapter]}
\floatname{codelisting}{Listing}
\newcommand*\listoflistings{\listof{codelisting}{List of Listings}}
\makeatother
\makeatletter
\makeatother
\makeatletter
\@ifpackageloaded{caption}{}{\usepackage{caption}}
\@ifpackageloaded{subcaption}{}{\usepackage{subcaption}}
\makeatother
\usepackage{bookmark}
\IfFileExists{xurl.sty}{\usepackage{xurl}}{} % add URL line breaks if available
\urlstyle{same}
\hypersetup{
  pdftitle={The Importance of Structure},
  pdfauthor={Olivier V. Raven; Robin Holmes; Ian A. K. Kusabs; Francis J. Burdon; Deniz Özkundakci},
  pdfkeywords={Paranephrops planifrons, Artificial
reef, Behaviour, Habitat preference, Aotearoa New Zealand},
  colorlinks=true,
  linkcolor={blue},
  filecolor={Maroon},
  citecolor={Blue},
  urlcolor={Blue},
  pdfcreator={LaTeX via pandoc}}


\title{The Importance of Structure}
\usepackage{etoolbox}
\makeatletter
\providecommand{\subtitle}[1]{% add subtitle to \maketitle
  \apptocmd{\@title}{\par {\large #1 \par}}{}{}
}
\makeatother
\subtitle{The importance of structure: mesocosm evidence for shelter
design in freshwater crayfish restoration}
\author{Olivier V. Raven \and Robin Holmes \and Ian A. K.
Kusabs \and Francis J. Burdon \and Deniz Özkundakci}
\date{2026-08-04}
\begin{document}
\maketitle


\clearpage
\providecommand{\chaptermark}[1]{}
\ifcsname c@chapter\endcsname
\setcounter{chapter}{4}
\setcounter{section}{0}
\addcontentsline{toc}{chapter}{\protect\numberline{4}The Importance of Structure}
\chaptermark{The Importance of Structure}
\fi
\providecommand{\thechapter}{4}
\thispagestyle{empty}
\begin{tikzpicture}[remember picture, overlay]
  \begin{scope}
    \clip ([yshift=-7cm]current page.north west)
          rectangle ([yshift=4.5cm]current page.south east);
    \node[inner sep=0pt, anchor=north]
      at ([yshift=-7cm]current page.north)
      {\includegraphics[width=18.2cm, height=\paperwidth, angle=90]{images/ch4_cover.jpg}};
  \end{scope}
  \fill[black]
    (current page.north west) rectangle ([yshift=-7cm]current page.north east);
  \node[color=thesissand, align=left,
        font=\fontsize{13}{16}\selectfont\sffamily, anchor=west]
    at ([xshift=1.5cm, yshift=-1.5cm]current page.north west)
    {Chapter \thechapter};
  \node[color=white, text width=0.85\paperwidth, align=left,
        font=\fontsize{22}{28}\selectfont\bfseries\sffamily, anchor=west]
    at ([xshift=1.5cm, yshift=-4cm]current page.north west)
    {The Importance of Structure};
  \node[color=white, text width=0.85\paperwidth, align=left,
        font=\fontsize{10}{13}\selectfont\itshape\sffamily, anchor=north west]
    at ([xshift=1.5cm, yshift=-5.0cm]current page.north west)
    {The importance of structure: mesocosm evidence for shelter design in freshwater crayfish restoration};
  \fill[black]
    (current page.south west) rectangle ([yshift=4.5cm]current page.south east);
  \node[color=white, text width=0.85\paperwidth, align=left,
        font=\fontsize{9}{13}\selectfont\sffamily, anchor=west]
    at ([xshift=1.5cm, yshift=3.2cm]current page.south west)
    {Olivier V.\ Raven, Robin Holmes, Ian A.\ K.\ Kusabs, Francis J.\ Burdon \& Deniz Özkundakci};
  \node[color=thesissand, text width=0.85\paperwidth, align=left,
        font=\fontsize{9}{13}\selectfont\itshape\sffamily, anchor=west]
    at ([xshift=1.5cm, yshift=1.5cm]current page.south west)
    {Aquatic Conservation: Marine and Freshwater Ecosystems --- Under review};
\end{tikzpicture}
\null
\clearpage

\subsection*{Abstract}\label{abstract}
\addcontentsline{toc}{subsection}{Abstract}

Freshwater crayfish are keystone species that rank among the most
threatened invertebrate groups globally. While habitat loss and
predation by non-native predators feature prominently as drivers of
crayfish population declines, the optimal design of habitat enhancement
structures for their restoration remains largely unknown. To better
understand freshwater crayfish habitat preferences, our study focused on
the endemic kōura (\emph{Paranephrops planifrons}) found in Aotearoa New
Zealand. We used four mesocosm experiments to examine how shelter
structure, material composition, individual body size, and group
dynamics influence refuge selection in kōura. Individual kōura strongly
preferred a wooden log with enclosed drilled holes, but when housed in
mixed-size groups, their preference shifted to flat and piled stone
configurations that provided more complex refugia for individuals of
different sizes. Flat stone structures were preferred over split wood
pieces of similar dimensions, and this was strongest for large
individuals. When structural dimensions were held constant, no
significant preference was detected between concrete and wooden
structures, indicating that material composition alone does not drive
refuge selection. The findings of these mesocosm experiments suggest
that stone pile configurations warrant field evaluations as habitat
enhancement structures for kōura restoration. The stone piles create
structural complexity that can provide habitat for different size
classes and thus may add critical refuge from non-native predators for
kōura populations.

\subsection{Introduction}\label{introduction}

Freshwater biodiversity is declining faster than in any other ecosystem
globally, driven by habitat loss, pollution, and invasive species
(\citeproc{ref-Dudgeon2006}{Dudgeon \emph{et al.}, 2006}). Crayfish are
among the most threatened freshwater invertebrates, with 32 \% of
species assessed as threatened with extinction due to these same
pressures (\citeproc{ref-Richman2015}{Richman \emph{et al.}, 2015}). As
keystone species and ecosystem engineers, crayfish play crucial roles in
freshwater food webs through shredding of leaf litter, bioturbation of
sediments, and nutrient cycling, and their loss can have profound
cascading effects on freshwater ecosystems
(\citeproc{ref-Momot1995}{Momot, 1995};
\citeproc{ref-Collier1997}{Collier \emph{et al.}, 1997};
\citeproc{ref-Parkyn2001}{Parkyn \emph{et al.}, 2001}). The increasing
degradation of freshwater habitats therefore not only threatens crayfish
populations but also risks disrupting the ecosystem processes crayfish
sustain.

Habitat complexity is a key driver of crayfish population persistence,
as structurally diverse habitats provide physical refugia from
predators, foraging opportunities, and shelter during moulting when
individuals are particularly vulnerable
(\citeproc{ref-Hammond2006}{Hammond \emph{et al.}, 2006}). Shelter
availability directly increases crayfish survival and growth under
predation pressure (\citeproc{ref-Adams2007}{Adams, 2007}), making
refuge quality a critical consideration for both natural habitat
persistence and restoration design. Crayfish use a wide variety of
natural and artificial structures as shelter, including rocks, woody
debris, and even anthropogenic objects such as bottles and discarded
appliances (\citeproc{ref-Parkyn2009}{Parkyn \emph{et al.}, 2009};
\citeproc{ref-Adams2014}{Adams, 2014}), demonstrating considerable
flexibility in habitat use. While this diversity of potential refugia is
well documented, preferences among different structure types remain
poorly understood, even though knowledge of habitat preferences is
fundamental to designing effective enhancement structures for
restoration. In the context of freshwater restoration, increasing
structural complexity through providing shelters has improved crayfish
abundance and restocking success in degraded streams and lakes
(\citeproc{ref-Huolila1997}{Huolila \emph{et al.}, 1997};
\citeproc{ref-Johnsen2008}{Johnsen and Taugbøl, 2008}), and the
availability of coarser substrate has also been shown to improve
juvenile survival rates (\citeproc{ref-Mazlum2017}{Mazlum \emph{et al.},
2017}). These findings indicate that enhancing habitat complexity can be
a useful management tool for threatened crayfish populations, yet
systematic comparisons of shelter configurations for freshwater crayfish
restoration remain largely absent from the literature.

Shelter selection in crayfish is strongly influenced by body size, as
individuals actively select refugia proportional to their dimensions.
Crayfish individuals tend to select stones more than 3 times longer and
1.25 times wider than their carapace length, with larger males
consistently selecting larger stones
(\citeproc{ref-Streissl2002}{Streissl and Hödl, 2002}). Social dynamics
further complicate shelter use when refuge availability is limited, as
competitive interactions among conspecifics determine which individuals
gain access to preferred shelters. Dominance hierarchies based on body
size and reproductive status show that males and maternal females often
are able to take over shelters from non-maternal females, while adults
consistently outcompete juveniles for shelter access
(\citeproc{ref-Figler2005}{Figler \emph{et al.}, 2005}). Body size and
shelter dimensions have direct competitive consequences as crayfish in
undersized shelters initiated and won more fights than those in
appropriately sized shelters, suggesting that shelter fit influences an
individual's assessment of its own fighting ability
(\citeproc{ref-Percival2010}{Percival and Moore, 2010}). Body size
ratios between contestants are a strong predictor of shelter takeover
success, although prior ownership can offset size disadvantages
(\citeproc{ref-Ranta1993}{Ranta and Lindström, 1993}), particularly when
both individuals hold high-quality shelters
(\citeproc{ref-Takahashi2019}{Takahashi \emph{et al.}, 2019}). At the
population level, habitat complexity buffers these competitive effects,
with juveniles in structurally complex habitats showing higher survival,
faster growth, and greater moulting rates as more shelter availability
reduces agonistic interactions (\citeproc{ref-Olsson2009}{Olsson and
Nyström, 2009}). Together, shelter use in crayfish reflects an
interaction between individual body size and social competitive
dynamics, suggesting that effective habitat enhancement structures must
provide a range of shelter sizes to support different life stages and
size classes simultaneously.

The material and structural properties of shelters can influence their
value as refugia. In a study comparing shelter configurations on reef
fish assemblages, the availability and size of enclosed spaces mattered
more than total structure volume (\citeproc{ref-Hylkema2020}{Hylkema
\emph{et al.}, 2020}). Additionally, enclosed structures such as PVC
pipes provide greater protection against predatory fish than stone or
vegetation structures, likely due to their consistent hole dimensions
(\citeproc{ref-Zhao2024}{Zhao \emph{et al.}, 2024}). However, the use of
synthetic materials like PVC pipes in restoration ecology brings
ecological concerns, particularly in contexts where introducing
non-native materials may conflict with conservation values or indigenous
cultural frameworks (\citeproc{ref-Wehi2017}{Wehi and Lord, 2017};
\citeproc{ref-Cooke2023}{Cooke \emph{et al.}, 2023}). Natural materials
such as stones and wood offer ecologically compatible alternatives, yet
the importance of structural configuration versus material composition
for crayfish shelter preference is poorly understood.

In Aotearoa-New Zealand (hereafter Aotearoa), the endemic freshwater
crayfish kōura (\emph{Paranephrops planifrons}) faces conservation
pressure from habitat loss, water quality degradation, and predation by
invasive species (\citeproc{ref-Lee2025}{Lee \emph{et al.}, 2025};
\citeproc{ref-Kusabs2026}{Kusabs \emph{et al.}, 2026}). Kōura are a
taonga (culturally treasured) species for Māori, the indigenous people
of Aotearoa, and hold significant ecological and cultural value in
freshwater ecosystems (\citeproc{ref-Kusabs2015a}{Kusabs \emph{et al.},
2015a}; \citeproc{ref-Kusabs2009}{Kusabs and Quinn, 2009}). In their
natural habitat, kōura are positively associated with overhanging banks,
coarse substrates, and woody debris (\citeproc{ref-Raven2026a}{Raven
\emph{et al.}, 2026a}; \citeproc{ref-Usio2000}{Usio and Townsend, 2000};
\citeproc{ref-Jowett2008}{Jowett \emph{et al.}, 2008};
\citeproc{ref-Parkyn2009}{Parkyn \emph{et al.}, 2009}). Kōura have been
shown to significantly increase refuge use in the presence of predators
(\citeproc{ref-Raven2026b}{Raven \emph{et al.}, 2026b};
\citeproc{ref-Shave1994}{Shave \emph{et al.}, 1994}). Habitat
restoration by increasing shelter availability represents a promising
management tool for kōura recovery, and the use of natural materials
aligns with ecological and cultural values.

In this study, four mesocosm experiments were conducted to identify
kōura habitat preferences. We assessed how different structures varying
in material and configuration influence refuge selection, accounting for
individual behaviour, body size, and group interactions. We predicted
that 1) Structurally complex structures providing enclosed crevices
would be preferred over isolated or open structures, as enclosed spaces
reduce predation risk; 2) structure preference would be size-dependent,
with larger kōura selecting structures with larger refuge spaces
(\citeproc{ref-Streissl2002}{Streissl and Hödl, 2002};
\citeproc{ref-Percival2010}{Percival and Moore, 2010}); 3) stone-based
structures would be preferred over wood due to greater structural
stability, and 4) material type alone would not influence refuge
preference between concrete and wooden bricks. The results provide
practical guidance for designing cost-effective, ecologically
appropriate habitat enhancement structures to enhance kōura abundance
and support population recovery.

\subsection{Materials and Methods}\label{materials-and-methods}

\subsubsection{Animal collection and tank
setup}\label{animal-collection-and-tank-setup}

In mid-August 2025, 72 kōura were sourced from Lakes Ōkāreka and
Tikitapu, located within the Rotorua Te Arawa Lakes area in the North
Island of Aotearoa. Collected animals were transported to the Aquatic
Research Centre at the University of Waikato in Hamilton. Individuals
were assigned to one of three size classes based on orbital carapace
length (OCL): small (\textless{} 22 mm), medium (22--30 mm), and large
(\textgreater{} 30 mm). Each animal was marked with a unique
identification number written on the carapace using an orange oil-based
paint marker (Sharpie, medium point) (\citeproc{ref-Ramalho2010}{Ramalho
\emph{et al.}, 2010}).

Kōura were acclimatised for 12 days in outdoor concrete holding tanks (Ø
1.4 m), each filled to 0.4 m with dechlorinated tap water conditioned
for three months, paving sand as substrate, shelters, and were covered
with a single layer of shade cloth. Each holding tank housed 10--12
individuals. Kōura were fed every second day with commercial carnivore
pellets (Hikari Tropical, Kyorin Food Industries Ltd., Japan; animal
protein 47 \%). Water temperature and dissolved oxygen remained stable
throughout the experimental period (temperature: 12 ± 1.2 °C; dissolved
oxygen: 99.9 ± 1.7 \%).

Experiments were conducted in six larger outdoor concrete tanks (1425 L;
Ø 1.85 m) with the same substrate and water conditions, covered with a
double layer of shade cloth to prevent escape and minimise solar
heating. Fresh dechlorinated tap water and an air stone were added
between rounds. To limit thermal stress, experiments were conducted
during the Australasian spring (September to October). Because kōura are
a taonga (treasured) species, permission to conduct the experiments was
obtained from Te Arawa Lakes Trust and Te Komiti Whakahaere. All
experimental procedures were conducted under University of Waikato
animal ethics committee approval (protocol number: 1207; see Ethical
Approval section).

\subsubsection{Experiment 1: Individual
preference}\label{experiment-1-individual-preference}

This experiment assessed individual kōura structure preference by
offering single animals a choice among four structures
(Figure~\ref{fig-exp-setup}). The structure-types were: SingleStone (ten
quarried stones (150--250 mm) placed individually across the quadrant
with sufficient spacing between stones that no crevices formed between
them), FlatStone (ten quarried stones arranged tightly without
stacking), StonePile (ten quarried stones stacked to form a pile), and
WoodLog (a tānekaha (\emph{Phyllocladus trichomanoides}) log (500 mm
long, 150 mm diameter) containing six drilled holes (70 mm deep, 32 mm
diameter)). If a kōura was not using a structure, TankWall was noted as
its location which represents the wall of the tank or when the kōura was
in the open. TankWall functioned as an indicator of structure rejection
rather than a preference. Six experimental tanks were used, each housing
a single medium-sized kōura. Across five experimental replicates, a
total of 30 kōura were tested, consisting of 15 females with mean OCL
and weight ± SD: 26.9 ± 2 mm, range: 23.2--29.8 mm, 12.5 ± 3.3 g and 15
males: 26.1 ± 2.2 mm, range: 22.6--30 mm, 11.9 ± 3.1 g. Each replicate
lasted 22 hours, beginning at noon (12:00) and ending at 10:00 the
following morning.

Each experimental tank was divided into four quadrants, with one
structure type placed at random (randomizer.org) in each quadrant. For
each trial, a single kōura was placed inside a vertically oriented
plastic tube (50 mm long, open at one end and sealed at the other)
positioned at the centre of the tank, with the sealed end resting on the
substrate. This ensured a standardised starting position from which the
kōura could exit in any direction (\citeproc{ref-MacNeil2025}{MacNeil
\emph{et al.}, 2025}). The time taken for the kōura to locate its
initial refuge location was recorded. On the following day, the final
location of the kōura was noted after which it was removed from the
tank. Tanks were reset and structure placements re-randomised prior to
the next trial. Each individual kōura was used once in a single trial
and not reused across temporal replicate blocks. To confirm nocturnal
activity, three digital cameras (GoPro Hero 3) were placed in one tank
and recorded for five hours at the end of each of the five days.

\subsubsection{Experiment 2: Group
dynamics}\label{experiment-2-group-dynamics}

Building on Experiment 1, this experiment examined how group dynamics
and body size influenced structure preference when kōura were housed
together. The same six experimental tanks were used and configured with
the identical structure types and randomisation procedure as in
Experiment 1 (Figure~\ref{fig-exp-setup}). In contrast to the
single-animal trials, each tank housed a group of six male kōura, as
male-only groups simplified competitive dynamics and size class
assignment given limited female availability. Groups were comprised of
two small, two medium, and two large individuals. Size classes were
defined by OCL: small (mean OCL and weight ± SD: 19.2 ± 2 mm, range:
15.8--21.6 mm, 3.9 ± 2.1 g), medium (26.4 ± 1.6 mm, 23.5--29.5 mm, 11.6
± 2.4 g), and large (36.3 ± 4.2 mm, 30.7--43.3 mm, 31.8 ± 11.2 g). Each
group was introduced to the centre of the tank within an open-top
plastic container (150 × 150 × 85 mm). The container was removed, and
the initial location of each individual was recorded. On the following
day, the final location of each kōura was documented, before all
individuals were removed from the tank. This procedure was repeated
across five experimental replicate blocks, resulting in 30 group-trials
in total (six groups × five blocks). Structure placements and tank
assignments were re-randomised between blocks, such that each group was
housed in a different tank in each successive round. The same six groups
were reused across all five blocks, with individuals staying with their
assigned group between replicates. No mortality occurred during the
experiment.

\subsubsection{Experiment 3: Wood vs
Stone}\label{experiment-3-wood-vs-stone}

Using the same male kōura groups from Experiment 2, this experiment
narrowed the choice to two structure types, stone and wood, to directly
compare preference between these materials. The structure types were:
(1) a flat stone structure (FlatStone) constructed from twelve stones
placed tightly together, and (2) a wood structure consisting of twelve
wood splits (WoodSplit) made from the tānekaha logs
(Figure~\ref{fig-exp-setup}). Each group was released at the centre of
the tank, and the initial structure selected by each kōura was recorded.
The following day, the final location of each kōura was documented. This
procedure was repeated across three experimental replicate blocks,
resulting in 18 group-trials in total (six groups × three blocks).
Structure placements and tank assignments were re-randomised between
blocks, such that each group was housed in a different tank in each
successive round. One small individual died in the first block and was
replaced with a size-matched individual to maintain group composition.
One individual could not be located at the end of the final block and
was assumed to have escaped the tank.

\subsubsection{Experiment 4: Material
comparison}\label{experiment-4-material-comparison}

This experiment tested whether kōura showed a preference between
perforated concrete bricks and perforated wooden blocks to further
identify the effect of material on structure preference. Two perforated
building bricks (ConcreteBrick) of 80 × 80 × 225 mm with five holes (Ø
30 mm) were tested against two wooden blocks (WoodBlock) with drilled
holes made from Western red cedar (\emph{Thuja plicata}), constructed to
match the building bricks in external dimensions and hole size. To
prevent floating, a solid clay brick of identical dimensions was placed
on top of each WoodBlock. ConcreteBrick did not require additional
weight (Figure~\ref{fig-exp-setup}). New groups of ten kōura and four
experiment tanks were used for this experiment, comprising a total of 40
individuals consisting of 18 females with mean OCL and weight ± SD: 26.2
± 3.1 mm, range: 20.2--30.8 mm, 12.2 ± 4.2 g and 22 males: 26 ± 2.1 mm,
range: 20.9--29.5 mm, 11.6 ± 3.2 g. Due to reduced animal availability,
individuals were drawn from a slightly broader size range than
medium-size individuals, resulting in unbalanced size class
representation (small: n = 4, medium: n = 35, large: n = 1). In each of
the four experimental tank quadrants one structure was placed, with
similar structure types positioned opposite to one another. Each group
of kōura was released at the centre of the tank and the following day,
the final location of each kōura was documented. This experiment was
repeated across four temporal blocks using the same four groups of ten
kōura, resulting in 16 group-trials in total (four groups × four
blocks). Structure placements and tank assignments were re-randomised
between blocks, such that each group was housed in a different tank in
each successive round.

\subsubsection{Data analysis}\label{data-analysis}

Experiment 1: To test the difference between the four structure types
selected by the single kōura, a multinomial goodness-of-fit test was
performed using the XNomial package (\citeproc{ref-Engels2014}{Engels,
2014}). TankWall was excluded from preference tests as it represents the
absence of structure choice rather than a structure type preference.
Additionally, pairwise likelihood-ratio G-tests with Bonferroni
correction were used to identify which structure types differed
significantly from each other. Bayesian categorical mixed-effects models
were used to test whether body size or sex affected structure preference
across all experiments, fitted using the brms package
(\citeproc{ref-Burkner2025}{Bürkner \emph{et al.}, 2025}). For
Experiments 1 and 2, where kōura chose among four structure types, a
categorical family was used. For Experiments 3 and 4, where only two
structure types were directly compared, a Bernoulli family was used.
Model convergence was assessed using the Gelman-Rubin diagnostic (all
\(\hat{R}\) = 1.00) and divergent transition counts (all experiments: 0
divergences), confirming reliable posterior sampling
(\citeproc{ref-Gelman1992}{Gelman and Rubin, 1992}). To test whether
structure type use differed between sexes, the Bayesian categorical
mixed-effects model was fitted with structure type as the outcome, sex
as a fixed effect and tank and experimental round as random intercepts.
A permutation-based symmetry test was applied to assess whether
structure type transitions between initial location and next-day
location were directional rather than random.

Experiment 2: To test whether overall structure preference existed at
the group level, one-sample Wilcoxon signed-rank tests were applied to
group-round proportions for each structure type, testing against a null
expectation of 0.25 (equal distribution among four structure types).
P-values were Bonferroni adjusted for multiple comparisons. Normal
approximation was used due to ties arising from discrete group counts.
To test whether size classes affected structure preference while
accounting for repeated measurement of individuals and groups, the
Bayesian categorical mixed-effects model was fitted with structure type
as the outcome, size class as a fixed effect, and group identity,
individual identity nested within group, and tank and experimental round
as random intercepts. The Stuart-Maxwell test was used to test for
statistical differences between initial and next-day location.

Experiment 3: To test the difference between FlatStone and WoodSplit the
Wilcoxon signed-rank test was applied to group-round proportions for
each structure type, testing against a null expectation of 0.5. P-values
were Bonferroni adjusted. To test whether size classes affected
structure preference while accounting for repeated measurement of
individuals and groups, the Bayesian Bernoulli mixed-effects model was
fitted with size class as a fixed effect, following the same random
effects structure as described for Experiment 2. The Stuart-Maxwell test
was used to test for statistical differences between initial and
next-day location.

Experiment 4: To test whether kōura showed a preference between
ConcreteBrick and WoodBlock, the Wilcoxon signed-rank test was applied
to group-round proportions for each structure type, testing against a
null expectation of 0.5 and p-values were Bonferroni adjusted. A paired
Wilcoxon signed-rank test was used to directly compare group-round
proportions between ConcreteBrick and WoodBlock. As size class was
unbalanced, the effect of sex on structure preference was tested using
the Bayesian Bernoulli mixed-effects model with sex as a fixed effect,
following the same random effects structure as described for Experiment
2. Finally, as only next-day locations were recorded in this experiment,
no directional shift analysis was performed.

\subsection{Results}\label{results}

\subsubsection{Experiment 1: Individual
preference}\label{experiment-1-individual-preference-1}

Individual kōura showed a strong preference for WoodLog, selecting it
significantly more often than SingleStone, FlatStone, and StonePile
(pairwise G-test, Bonferroni adjusted: \emph{p} \textless{} 0.001,
0.011, 0.004, respectively; overall multinomial test: \emph{p}
\textless{} 0.001; Figure~\ref{fig-exp1}). No significant differences
were found among the stone structure types. Overall, WoodLog was
selected by 67 \% of kōura at refuge (males: 80 \%; females: 53 \%), and
no kōura were found at TankWall. Sex had no credible effect on structure
preference (β = 3.32, 95\% CI: -1.09 to 9.3; posterior probability =
0.92) with credible intervals spanning zero for all structure types. The
symmetry test indicated that structure type use shifted between initial
and next-day location (\emph{p} = 0.001), suggesting kōura actively
relocated to preferred structure types overnight rather than remaining
at their initial location. The video footage showed kōura emerging,
walking and climbing on the structures after sunset.

\subsubsection{Experiment 2: Group
dynamics}\label{experiment-2-group-dynamics-1}

When housed in mixed-size groups, kōura showed a significant overall
preference for stone-based structures (Figure~\ref{fig-exp2}). FlatStone
was used significantly more (\emph{p} = 0.004), while SingleStone
(\emph{p} \textless{} 0.001) and WoodLog (\emph{p} = 0.039) were
significantly avoided. StonePile did not differ from random expectation
(\emph{p} = 1.000) and did not differ significantly from FlatStone
(\emph{p} = 0.117). No kōura were found at TankWall, indicating active
structure selection across the group. Size classes had limited but
credible effect on structure preference. Medium-sized kōura showed a
weaker preference for FlatStone over SingleStone compared to small kōura
(β = -1.77, 95\% CI: -3.7 to -0.14; posterior probability = 0.99), while
no other size class effect was credible. Overall, FlatStone remained the
most selected structure type across all size classes, but its strength
of preference declined from small (posterior probability of top choice =
82 \%) to medium (53 \%), and with large kōura showing similar
preference between FlatStone and StonePile (42 \% vs 38 \%). Structure
type use shifted directionally between initial and next-day location
(Stuart-Maxwell test: \emph{p} = 0.001), indicating that kōura actively
relocated to preferred structure types overnight.

\subsubsection{Experiment 3: Wood vs
Stone}\label{experiment-3-wood-vs-stone-1}

Kōura consistently selected FlatStone over WoodSplit, with the
preference growing with body size (Figure~\ref{fig-exp3}). At the group
level, FlatStone was used significantly more than expected (\emph{p} =
0.001), while WoodSplit was significantly avoided (\emph{p} = 0.004).
Only one kōura was found at TankWall. The Bayesian model indicated that
preference for FlatStone strengthened with body size, with the posterior
probability of FlatStone being the top-ranked structure increasing from
small (66 \%) to medium (79 \%), to large kōura (99 \%). Large kōura
showed a particularly strong and credible preference for FlatStone over
WoodSplit (β = -2.55, 95\% CI: -4.7 to -0.9; posterior probability =
\textgreater{} 0.99), while medium kōura did not differ from small
kōura. Structure type use shifted directionally between initial and
next-day location (Stuart-Maxwell test: \emph{p} = 0.003), with kōura
relocating toward FlatStone overnight.

\subsubsection{Experiment 4: Material
comparison}\label{experiment-4-material-comparison-1}

Kōura used both structure types, showing no significant preference
between ConcreteBrick and WoodBlock (\emph{p} = 0.137;
Figure~\ref{fig-exp4}), with mean group proportions of 47 \% and 53 \%
respectively. Neither structure type differed significantly from equal
distribution (\emph{p} = 0.273). 9 \% of kōura were found at TankWall,
which is higher than in any previous experiment, likely reflecting the
smaller total structure size in this experiment rather than active
structure avoidance. Sex had no credible effect on preference between
structure types (β = -0.18, 95\% CI: -1.24 to 0.84; posterior
probability = 0.65), indicating that both structures were equally
suitable refuge for male and female kōura.

\subsection{Discussion}\label{discussion}

\subsubsection{Structure matters more than
material}\label{structure-matters-more-than-material}

The consistent finding across all experiments that structural stability
and cavities complexity determined preference, regardless of whether the
comparison involved material type, configuration or group size, suggests
that kōura, like many refuge-dependent crustaceans
(\citeproc{ref-ZimmerFaust1987}{Zimmer-Faust and Spanier, 1987};
\citeproc{ref-Beck1995}{Beck, 1995};
\citeproc{ref-Streissl2002}{Streissl and Hödl, 2002};
\citeproc{ref-Pirtle2012}{Pirtle \emph{et al.}, 2012}), are primarily
selecting for physical shelter quality rather than material familiarity.
When hole dimensions were held constant between concrete and wooden
perforated structures, no preference for either material was detected,
demonstrating that material composition alone may not drive refuge
selection when structure is equivalent. Neither structure type was
conditioned prior to experiments, meaning biofilm colonisation that
occurs in natural lake environments was absent. As stone and wood
develop different biofilm communities (\citeproc{ref-Tank1996}{Tank and
Winterbourn, 1996}), material preference in field settings could change.
The wooden structures required a clay brick placed on top to prevent
floating, creating a difference with the concrete structures, though the
absence of any preference suggests this did not meaningfully influence
the result. This material-independence broke down when material affected
structural properties: when flat-laid stones were compared against
similarly sized split wood pieces, kōura showed a clear preference for
stone. The wood splits, being less dense than stone, sank but created
less structurally stable refuge places in the sandy substrate than the
stones, as larger kōura were able to move them {[}O.V. Raven, pers.
obs.{]}.

In group settings, structures providing the most crevice complexity were
strongly preferred, with flat-laid stones accommodating the most kōura,
followed by stone piles. Individual kōura showed a different preference,
consistently choosing the perforated wooden log over all stone options,
which is consistent with our first prediction that structurally complex
and enclosed refugia would be preferred as enclosed spaces reduce
predation risk (\citeproc{ref-Jurcak2016}{Jurcak \emph{et al.}, 2016}).
The individual drilled holes in the wooden log offered protection
similar to what has been found for perforated brick and PVC tube
shelters in response to the threat of catfish predation
(\citeproc{ref-Adams2007}{Adams, 2007}; \citeproc{ref-Zhao2024}{Zhao
\emph{et al.}, 2024}). The directional shift in structure type use,
between initial locations shortly after release and their location
following a night within the tank, confirmed across all experiments that
kōura actively relocated overnight, supporting the interpretation that
next-day locations reflect genuine refuge preference rather than chance
displacement.

The shift from the perforated wooden log in individual trials to stone
structures in group trials was unexpected. The wooden log contained six
individual drilled holes, so capacity was not limiting for the group of
six kōura. The preference shift may partly reflect size mismatch between
fixed hole dimensions, and the range of body sizes present in each
group, with only medium-sized individual finding the holes well-suited,
while small and large individuals were better served by the range of
crevices created by the stone structures
(\citeproc{ref-Streissl2002}{Streissl and Hödl, 2002}). However, a more
fundamental explanation may relate to social behaviour. The wooden log
was the most selected structure in round 1 of Experiment 2 (17 out of 36
times), similar to Experiment 1, but this preference did not persist in
subsequent rounds. This shift suggests that once kōura had experienced
sharing the tank, they adjusted their shelter use toward structures
allowing closer conspecific proximity. Gregarious shelter behaviour has
been documented in juvenile burrowing crayfish
(\citeproc{ref-Hazlett2007}{Hazlett \emph{et al.}, 2007};
\citeproc{ref-Reynolds2013}{Reynolds \emph{et al.}, 2013};
\citeproc{ref-Clay2017}{Clay \emph{et al.}, 2017}) and is well
established in lobsters including the New Zealand rock lobster
(\emph{Jasus edwardsii} Hutton, 1875) (\citeproc{ref-Butler1999}{Butler
\emph{et al.}, 1999}), where conspecific chemical cues attract
individuals to share refugia and find refugia significantly faster,
reducing predation exposure (\citeproc{ref-ZimmerFaust1987}{Zimmer-Faust
and Spanier, 1987}; \citeproc{ref-Childress2001}{Childress and
Herrnkind, 2001}). Whether kōura exhibit similar gregarious tendencies
remains understudied, but the preference shift observed in this study,
might suggest that social context alters shelter preferences.

\subsubsection{Body size and competitive
dynamics}\label{body-size-and-competitive-dynamics}

Size class consistently influenced refuge preference across experiments,
but the pattern differed by context. In Experiment 3, larger kōura
showed a stronger preference for flat stones over wooden splits than
smaller individuals, consistent with size dependent shelter selection
documented in crayfish (\citeproc{ref-Streissl2002}{Streissl and Hödl,
2002}; \citeproc{ref-Percival2010}{Percival and Moore, 2010}). This may
reflect larger kōura being more capable of asserting dominance over
preferred stable refugia through competitive displacement, which is
consistent with the size-based dominance hierarchy documented in kōura
populations (\citeproc{ref-Ranta1993}{Ranta and Lindström, 1993};
\citeproc{ref-Figler2005}{Figler \emph{et al.}, 2005};
\citeproc{ref-Stewart2011}{Stewart and Tabak, 2011}). In Experiment 2,
where four structure types were available, small kōura showed the
strongest preference for flat stones while large kōura distributed more
evenly between flat stones and stone piles, suggesting that the large
continuous surface area of the flat stones provides refuge spaces
proportional to smaller body dimensions
(\citeproc{ref-Streissl2002}{Streissl and Hödl, 2002}), while larger
individuals find adequate crevice space in either stone configuration.
The absence of sex effects in Experiments 1 and 4 suggests that refuge
preference in kōura is driven primarily by body size and competitive
dominance rather than sex (\citeproc{ref-Figler1999}{Figler \emph{et
al.}, 1999}). When refuge availability is limited, competitive
hierarchies displace smaller and recently moulted individuals from
preferred refugia, leaving them exposed (\citeproc{ref-Ranta1993}{Ranta
and Lindström, 1993}; \citeproc{ref-Figler2005}{Figler \emph{et al.},
2005}). Providing sufficient refuge spaces with a range of crevice size
can reduce agonistic interactions and benefit the whole population
(\citeproc{ref-Olsson2009}{Olsson and Nyström, 2009}), supporting
complex structure configurations over structures with fixed hole
dimensions.

\subsubsection{Restoration implications}\label{restoration-implications}

With the threat from invasive predatory fish, such as catfish preying on
kōura (\citeproc{ref-Kusabs2026}{Kusabs \emph{et al.}, 2026}), enhancing
habitat complexity may reduce top-down control on kōura populations.
Additional habitat as used here can provide shelter spaces and offers a
promising management tool, as refuge availability directly influences
predation survival (\citeproc{ref-Adams2007}{Adams, 2007};
\citeproc{ref-Zhao2024}{Zhao \emph{et al.}, 2024}). The present study
demonstrates that structural configuration and stability drive refuge
preference more strongly than material composition, with stone-based
configurations preferred by groups and no preference detected between
materials when dimensions were equal, suggesting that stone-based
habitat enhancement structures warrant field evaluation in lake
environments.

The broader restoration literature supports both the structural
principle and the choice of natural materials. Structurally stable
crevices are important for freshwater species like crayfish
(\citeproc{ref-Huolila1997}{Huolila \emph{et al.}, 1997};
\citeproc{ref-Streissl2002}{Streissl and Hödl, 2002}), but also for
marine environments, as reef fish, marine shrimp, and lobsters have all
been shown to benefit from habitat enhancement structures such as
artificial reefs (\citeproc{ref-Spanier1994}{Spanier, 1994};
\citeproc{ref-Magnelia2008}{Magnelia \emph{et al.}, 2008};
\citeproc{ref-Hylkema2020}{Hylkema \emph{et al.}, 2020};
\citeproc{ref-Dickson2023}{Dickson \emph{et al.}, 2023}). Structures
made from stone and wood have been shown to outperform synthetic
alternatives and bare substrates (\citeproc{ref-Huolila1997}{Huolila
\emph{et al.}, 1997}; \citeproc{ref-Magnelia2008}{Magnelia \emph{et
al.}, 2008}; \citeproc{ref-Cooke2023}{Cooke \emph{et al.}, 2023};
\citeproc{ref-Dickson2023}{Dickson \emph{et al.}, 2023}), and do not
contribute to plastic pollution in the environment
(\citeproc{ref-Cooke2023}{Cooke \emph{et al.}, 2023}). Locally sourced
natural materials also align with \emph{kaitiakitanga}, which is the
Māori principle of environmental guardianship and preserves the natural
character of culturally significant lake environments
(\citeproc{ref-Clapcott2018}{Clapcott \emph{et al.}, 2018}).
Furthermore, existing habitat surveys suggest that kōura are associated
with coarser substrates (\citeproc{ref-Kusabs2015b}{Kusabs \emph{et
al.}, 2015b}; \citeproc{ref-Olsson2006}{Olsson \emph{et al.}, 2006};
\citeproc{ref-Jowett2008}{Jowett \emph{et al.}, 2008}), and cobbles have
previously been used in stream and lake restoration
(\citeproc{ref-Huolila1997}{Huolila \emph{et al.}, 1997};
\citeproc{ref-Johnsen2008}{Johnsen and Taugbøl, 2008}). For in-lake
deployment stone pile structures are likely to be the most practical
configuration, due to their elevated structure, stone piles are less
likely to be buried into soft substrates than flat individual rocks
(\citeproc{ref-Garron2024}{Garron \emph{et al.}, 2024}). Field
validation in lake environments remains an essential next step, as
in-lake conditions including water movement, natural substrate, and
predator pressure may modify kōura preferences in ways not captured
under controlled conditions.

\subsubsection{Conclusion}\label{conclusion}

Across four mesocosm experiments, structural configuration and crevice
complexity consistently predicted kōura shelter preference more strongly
than material composition. The shift from individual preference for
enclosed wooden log shelters to group preference for stone
configurations shows that social context plays a key role in determining
shelter value for kōura. Size-dependent preferences further suggest that
structures accommodating multiple body sizes simultaneously benefit
whole populations rather than only dominant individuals. Stone-based
structures using locally sourced materials are promising candidates for
field evaluation, aligning with ecological and cultural values
associated with restoring this taonga species. Field validation remains
the essential next step for translating these findings into effective
kōura habitat restoration practice.

\subsection{Acknowledgements}\label{acknowledgements}

We thank Te Arawa Lakes Trust and Te Komiti Whakahaere for the
opportunity to collect and study the treasured kōura. We are grateful to
Soweeta Fort-D'ath and William Anaru (Te Arawa Lakes Trust), and Tihini
Grant (Ngāti Pikiao).

\subsubsection{Funding statement}\label{funding-statement}

This research was supported by the Fish Futures programme funded through
a Ministry of Business, Innovation and Employment grant (CAWX2101) with
additional funding provided by the Bay of Plenty Regional Council under
the Toihuarewa Waimāori - Bay of Plenty Regional Council Chair in Lake
and Freshwater Science programme.

\subsubsection{Authors' contributions}\label{authors-contributions}

Study conception and design were undertaken by O.V.R., R.H. and F.J.B.
Experiments were conducted by O.V.R. Data analysis was performed by
O.V.R. I.A.K.K. obtained iwi consent, provided local ecological
knowledge and kōura specimens. Manuscript drafting was led by O.V.R. All
authors contributed to manuscript revision, approved the results, and
the final manuscript.

\subsubsection{Ethical approval}\label{ethical-approval}

This study received ethical approval by University of Waikato animal
ethics committee, protocol number: 1207. Additional permission to
conduct experiments with kōura was granted by Te Arawa Lakes Trust and
Te Komiti Whakahaere.

\subsubsection{Declarations}\label{declarations}

The corresponding author has declared that none of the authors has any
competing interests.

\subsubsection{Data availability
statement}\label{data-availability-statement}

The code and derived data supporting the findings of this study are
openly available on GitHub
(https://github.com/OlivierRaven/Koura\_Mesocosm)and a rendered version
of the analysis is hosted at
https://olivierraven.github.io/Koura\_Mesocosm/. The primary raw dataset
is archived on Zenodo (https://doi.org/\ldots). Peer review documents,
including the cover letter and reviewer responses, are available in the
manuscript repository in the interest of open and transparent science.

\subsection{Figures}\label{figures}

\begin{figure}

\centering{

\pandocbounded{\includegraphics[keepaspectratio]{images/Setup_exp_tanks.png}}

}

\caption{\label{fig-exp-setup}Overhead photographs of experimental tank
configurations for each of the four mesocosm experiments. Experiment 1
(Individual preference) and Experiment 2 (Group dynamics): four shelter
types, WoodLog (tānekaha log with drilled holes), FlatStone (stones
arranged flat), StonePile (stones stacked), and SingleStone (individual
stones placed separately). Experiment 3 (Wood vs Stone): WoodSplit
shelter (tānekaha wood splits) and FlatStone shelter placed in opposing
halves of the tank. Experiment 4 (Material comparison): ConcreteBrick
(perforated concrete bricks) and WoodBlock (perforated wooden blocks of
identical dimensions) placed in opposing quadrants of the tank.}

\end{figure}%

\begin{figure}

\centering{

\includegraphics[width=5in,height=\textheight,keepaspectratio]{outputs/fig-exp1.png}

}

\caption{\label{fig-exp1}Proportion of individual kōura (Experiment 1)
found at each structure type the next day, shown separately for female
(white bars) and male (black bars) kōura. Bars show observed proportions
with sample sizes (n) shown on bars. Letters above bars indicate results
of pairwise G-tests with Bonferroni correction within each sex group,
structure types sharing the same letter within a panel do not differ
significantly from each other, while structure types with different
letters differ significantly in use. No significant difference in
overall structure preference was detected between sexes (see results).}

\end{figure}%

\begin{figure}

\centering{

\includegraphics[width=6.5in,height=\textheight,keepaspectratio]{outputs/fig-exp2.png}

}

\caption{\label{fig-exp2}Structure type preference of group kōura
(Experiment 2). (A) Proportion of kōura found at each structure type at
L1 across all groups and rounds. The dashed line indicates the null
expectation of equal distribution (0.25). Significance symbols indicate
deviation from null expectation (Wilcoxon signed-rank test, Bonferroni
adjusted): \(\uparrow\) used more than expected, \(\downarrow\) used
less than expected. Sample sizes shown in white. (B) Posterior predicted
probability of structure type use by size class from Bayesian
categorical mixed model, with 95 \% credible intervals.}

\end{figure}%

\begin{figure}

\centering{

\includegraphics[width=6.5in,height=\textheight,keepaspectratio]{outputs/fig-exp3.png}

}

\caption{\label{fig-exp3}Structure type preference in stone versus wood
comparison (Experiment 3). (A) Proportion of kōura found at each
structure type at next-day location across all groups and rounds. The
dashed line indicates the null expectation of equal distribution (0.5).
Significance symbols as in Figure 3. (B) Posterior predicted probability
of structure type use by size class from Bayesian categorical mixed
model, with 95 \% credible intervals. \(\uparrow\) used more than
expected, \(\downarrow\) used less than expected. TankWall excluded from
model predictions.}

\end{figure}%

\begin{figure}

\centering{

\includegraphics[width=6.5in,height=\textheight,keepaspectratio]{outputs/fig-exp4.png}

}

\caption{\label{fig-exp4}Material preference (Experiment 4). (A)
Proportion of kōura found at each structure type at next-day location
across all groups and rounds. The dashed line indicates the null
expectation of equal distribution (0.5). (B) Posterior predicted
probability of structure type use by sex from Bayesian categorical mixed
model, with 95 \% credible intervals. TankWall excluded from model
predictions.}

\end{figure}%

\protect\phantomsection\label{refs}
\begin{CSLReferences}{1}{0}
\bibitem[\citeproctext]{ref-Adams2007}
Adams SB. 2007. Direct and indirect effects of channel catfish
({\emph{Ictalurus punctatus}}) on native crayfishes ({Cambaridae}) in
experimental tanks. \emph{The American Midland Naturalist} \textbf{158}:
85--96.
https://doi.org/\href{https://doi.org/10.1674/0003-0031(2007)158\%5B85:daieoc\%5D2.0.co;2}{10.1674/0003-0031(2007)158{[}85:daieoc{]}2.0.co;2}

\bibitem[\citeproctext]{ref-Adams2014}
Adams SB. 2014. Crayfish use of trash versus natural cover in incised,
sand-bed streams. \emph{Environmental Management} \textbf{53}: 382--392.
https://doi.org/\href{https://doi.org/10.1007/s00267-013-0197-3}{10.1007/s00267-013-0197-3}

\bibitem[\citeproctext]{ref-Beck1995}
Beck MW. 1995. Size-specific shelter limitation in stone crabs: A test
of the demographic bottleneck hypothesis. \emph{Ecology} \textbf{76}:
968--980.
https://doi.org/\href{https://doi.org/10.2307/1939360}{10.2307/1939360}

\bibitem[\citeproctext]{ref-Burkner2025}
Bürkner P-C, Gabry J, Weber S, Johnson A. 2025.
\href{https://cran.r-project.org/web/packages/brms/index.html}{{brms}:
{Bayesian} regression models using {Stan}}.

\bibitem[\citeproctext]{ref-Butler1999}
Butler M, MacDiarmid A, Booth J. 1999. The cause and consequence of
ontogenetic changes in social aggregation in {New Zealand} spiny
lobsters. \emph{Marine Ecology Progress Series} \textbf{188}: 179--191.
https://doi.org/\href{https://doi.org/10.3354/meps188179}{10.3354/meps188179}

\bibitem[\citeproctext]{ref-Childress2001}
Childress MJ, Herrnkind WF. 2001. The guide effect influence on the
gregariousness of juvenile caribbean spiny lobsters. \emph{Animal
Behaviour} \textbf{62}: 465--472.
https://doi.org/\href{https://doi.org/10.1006/anbe.2001.1760}{10.1006/anbe.2001.1760}

\bibitem[\citeproctext]{ref-Clapcott2018}
Clapcott J, Ataria J, Hepburn C, Hikuroa D, Jackson A, Kirikiri R,
Williams E. 2018. {Mātauranga Māori}: Shaping marine and freshwater
futures. \emph{New Zealand Journal of Marine and Freshwater Research}
\textbf{52}: 457--466.
https://doi.org/\href{https://doi.org/10.1080/00288330.2018.1539404}{10.1080/00288330.2018.1539404}

\bibitem[\citeproctext]{ref-Clay2017}
Clay M, Stoeckel J, Helms B. 2017. The role of abiotic and biotic cues
in burrow habitat selection by juvenile crayfish. \emph{Behaviour}
\textbf{154}: 1177--1196.
https://doi.org/\href{https://doi.org/10.1163/1568539x-00003463}{10.1163/1568539x-00003463}

\bibitem[\citeproctext]{ref-Collier1997}
Collier KJ, Parkyn SM, Rabeni CF. 1997. Koura: A keystone species.
\emph{Water and Atmosphere} \textbf{5}: 18--20.

\bibitem[\citeproctext]{ref-Cooke2023}
Cooke SJ, Piczak ML, Vermaire JC, Kirkwood AE. 2023. On the troubling
use of plastic {`habitat'} structures for fish in freshwater ecosystems
-- or -- when restoration is just littering (D Lesbarrères, Ed).
\emph{FACETS} \textbf{8}: 1--19.
https://doi.org/\href{https://doi.org/10.1139/facets-2022-0210}{10.1139/facets-2022-0210}

\bibitem[\citeproctext]{ref-Dickson2023}
Dickson J, Franken O, Watson MS, Monnich B, Holthuijsen S, Eriksson BK,
Govers LL, Heide T van der, Bouma TJ. 2023. Who lives in a pear tree
under the sea? A first look at tree reefs as a complex natural
biodegradable structure to enhance biodiversity in marine systems.
\emph{Frontiers in Marine Science} \textbf{10}: 1213790.
https://doi.org/\href{https://doi.org/10.3389/fmars.2023.1213790}{10.3389/fmars.2023.1213790}

\bibitem[\citeproctext]{ref-Dudgeon2006}
Dudgeon D, Arthington AH, Gessner MO, Kawabata Z, Knowler DJ, Lévêque C,
Naiman RJ, Prieur‐Richard A, Soto D, Stiassny MLJ, et al. 2006.
Freshwater biodiversity: Importance, threats, status and conservation
challenges. \emph{Biological Reviews} \textbf{81}: 163--182.
https://doi.org/\href{https://doi.org/10.1017/s1464793105006950}{10.1017/s1464793105006950}

\bibitem[\citeproctext]{ref-Engels2014}
Engels B. 2014. {XNomial}: Exact goodness-of-fit test for multinomial
data with fixed probabilities. \emph{CRAN: Contributed Packages}.
https://doi.org/\href{https://doi.org/10.32614/cran.package.xnomial}{10.32614/cran.package.xnomial}

\bibitem[\citeproctext]{ref-Figler2005}
Figler MH, Blank GS, Peeke HVS. 2005. Shelter competition between
resident male red swamp crayfish {\emph{Procambarus clarkii}} (girard)
and conspecific intruders varying by sex and reproductive status.
\emph{Marine and Freshwater Behaviour and Physiology} \textbf{38}:
237--248.
https://doi.org/\href{https://doi.org/10.1080/10236240500376477}{10.1080/10236240500376477}

\bibitem[\citeproctext]{ref-Figler1999}
Figler MH, Cheverton HM, Blank GS. 1999. Shelter competition in juvenile
red swamp crayfish ({\emph{Procambarus clarkii}}): The influences of sex
differences, relative size, and prior residence. \emph{Aquaculture}
\textbf{178}: 63--75.
https://doi.org/\href{https://doi.org/10.1016/s0044-8486(99)00114-3}{10.1016/s0044-8486(99)00114-3}

\bibitem[\citeproctext]{ref-Garron2024}
Garron C, Pangaud A, Taubaty M, Rolan‐Meynard M, Argillier C, Raymond J,
Dublon J, Westrelin S. 2024. Towards the development of a standardized
method to monitor lakeshore hydromorphological restoration measures.
\emph{Aquatic Conservation: Marine and Freshwater Ecosystems}
\textbf{34}.
https://doi.org/\href{https://doi.org/10.1002/aqc.4152}{10.1002/aqc.4152}

\bibitem[\citeproctext]{ref-Gelman1992}
Gelman A, Rubin DB. 1992. Inference from iterative simulation using
multiple sequences. \emph{Statistical Science} \textbf{7}: 457--472.
https://doi.org/\href{https://doi.org/10.1214/ss/1177011136}{10.1214/ss/1177011136}

\bibitem[\citeproctext]{ref-Hammond2006}
Hammond KS, Hollows JW, Townsend CR, Lokman PM. 2006. Effects of
temperature and water calcium concentration on growth, survival and
moulting of freshwater crayfish, {\emph{Paranephrops zealandicus}}.
\emph{Aquaculture} \textbf{251}: 271--279.
https://doi.org/\href{https://doi.org/10.1016/j.aquaculture.2005.05.032}{10.1016/j.aquaculture.2005.05.032}

\bibitem[\citeproctext]{ref-Hazlett2007}
Hazlett BA, Lawler S, Edney G. 2007. Agonistic behavior of the crayfish
{\emph{Euastacus armatus}} and {\emph{Cherax destructor}}. \emph{Marine
and Freshwater Behaviour and Physiology} \textbf{40}: 257--266.
https://doi.org/\href{https://doi.org/10.1080/10236240701562412}{10.1080/10236240701562412}

\bibitem[\citeproctext]{ref-Huolila1997}
Huolila M, Marjomäki TJ, Laukkanen E. 1997. The success of crayfish
stocking in a dredged river with and without artificial shelter
increase. \emph{Fisheries Research} \textbf{32}: 185--189.
https://doi.org/\href{https://doi.org/10.1016/s0165-7836(97)00046-5}{10.1016/s0165-7836(97)00046-5}

\bibitem[\citeproctext]{ref-Hylkema2020}
Hylkema A, Debrot AO, Osinga R, Bron PS, Heesink DB, Izioka AK, Reid CB,
Rippen JC, Treibitz T, Yuval M, et al. 2020. Fish assemblages of three
common artificial reef designs during early colonization.
\emph{Ecological Engineering} \textbf{157}: 105994.
https://doi.org/\href{https://doi.org/10.1016/j.ecoleng.2020.105994}{10.1016/j.ecoleng.2020.105994}

\bibitem[\citeproctext]{ref-Johnsen2008}
Johnsen SI, Taugbøl T. 2008. Add stones, get crayfish -- is it that
simple? \emph{Freshwater Crayfish} \textbf{16}: 47--50.

\bibitem[\citeproctext]{ref-Jowett2008}
Jowett IG, Parkyn SM, Richardson J. 2008. Habitat characteristics of
crayfish ({\emph{Paranephrops planifrons}}) in {New Zealand} streams
using generalised additive models ({GAMs}). \emph{Hydrobiologia}
\textbf{596}: 353--365.
https://doi.org/\href{https://doi.org/10.1007/s10750-007-9108-z}{10.1007/s10750-007-9108-z}

\bibitem[\citeproctext]{ref-Jurcak2016}
Jurcak AM, Lahman SE, Wofford SJ, Moore PA. 2016. Behavior of crayfish.
In \emph{Biology and ecology of crayfish} CRC Press;

\bibitem[\citeproctext]{ref-Kusabs2015a}
Kusabs IA, Hicks BJ, Quinn JM, Hamilton DP. 2015a. Sustainable
management of freshwater crayfish (k{ō}ura, {\emph{Paranephrops
planifrons}}) in {Te Arawa} ({Rotorua}) lakes, {North Island, New
Zealand}. \emph{Fisheries Research} \textbf{168}: 35--46.
https://doi.org/\href{https://doi.org/10.1016/j.fishres.2015.03.015}{10.1016/j.fishres.2015.03.015}

\bibitem[\citeproctext]{ref-Kusabs2026}
Kusabs IAK, Duggan IC, Bruere AC, Scholes P, Te Kurapa TW, Grant TM,
Burdon FJ. 2026. Invasive catfish linked to the decline of freshwater
crayfish in two {Aotearoa New Zealand} lakes. \emph{Biological
Invasions} \textbf{28}.
https://doi.org/\href{https://doi.org/10.1007/s10530-026-03852-0}{10.1007/s10530-026-03852-0}

\bibitem[\citeproctext]{ref-Kusabs2009}
Kusabs IA, Quinn JM. 2009. Use of a traditional {Maori} harvesting
method, the tau kōura, for monitoring kōura (freshwater crayfish,
{\emph{Paranephrops planifrons}}) in {Lake Rotoiti}, {North Island, New
Zealand}. \emph{New Zealand Journal of Marine and Freshwater Research}
\textbf{43}: 713--722.
https://doi.org/\href{https://doi.org/10.1080/00288330909510036}{10.1080/00288330909510036}

\bibitem[\citeproctext]{ref-Kusabs2015b}
Kusabs IA, Quinn JM, Hamilton DP. 2015b. Effects of benthic substrate,
nutrient enrichment and predatory fish on freshwater crayfish (k{ō}ura,
{\emph{Paranephrops planifrons}}) population characteristics in seven
{Te Arawa} ({Rotorua}) lakes, {North Island, New Zealand}. \emph{Marine
and Freshwater Research} \textbf{66}: 631--643.
https://doi.org/\href{https://doi.org/10.1071/MF14148}{10.1071/MF14148}

\bibitem[\citeproctext]{ref-Lee2025}
Lee F, Kusabs IAK, Perry GLW, MacNeil C. 2025. Identifying refugia from
the synergistic threats of climate change and invasive species.
\emph{Web Ecology} \textbf{25}: 221--239.
https://doi.org/\href{https://doi.org/10.5194/we-25-221-2025}{10.5194/we-25-221-2025}

\bibitem[\citeproctext]{ref-MacNeil2025}
MacNeil C, Lee F, Kusabs I, Holmes R. 2025. The biter bit; scavenging
behaviour of native freshwater crayfish on the carrion of native and
introduced fish predators in {Aotearoa-New Zealand}. \emph{Knowledge \&
Management of Aquatic Ecosystems}: 29.
https://doi.org/\href{https://doi.org/10.1051/kmae/2025022}{10.1051/kmae/2025022}

\bibitem[\citeproctext]{ref-Magnelia2008}
Magnelia SJ, Jesus MJD, Schlechte JW, Cummings GC, Duty JL. 2008.
Comparison of plastic pipe and juniper tree fish attractors in a central
texas reservoir.

\bibitem[\citeproctext]{ref-Mazlum2017}
Mazlum Y, Guner Gurlek O, Sirin S. 2017. Effect of different substrates
on survival and growth of juveniles of the freshwater narrow-clawed
crayfish {\emph{Astacus leptodactylus}} eschscholtz, 1823 ({Decapoda,
Astacidae}). \emph{Crustaceana} \textbf{90}: 1289--1302.
https://doi.org/\href{https://doi.org/10.1163/15685403-00003724}{10.1163/15685403-00003724}

\bibitem[\citeproctext]{ref-Momot1995}
Momot WT. 1995. Redefining the role of crayfish in aquatic ecosystems.
\emph{Reviews in Fisheries Science} \textbf{3}: 33--63.
https://doi.org/\href{https://doi.org/10.1080/10641269509388566}{10.1080/10641269509388566}

\bibitem[\citeproctext]{ref-Olsson2009}
Olsson K, Nyström P. 2009. Non-interactive effects of habitat complexity
and adult crayfish on survival and growth of juvenile crayfish
({\emph{Pacifastacus leniusculus}}). \emph{Freshwater Biology}
\textbf{54}: 35--46.
https://doi.org/\href{https://doi.org/10.1111/j.1365-2427.2008.02089.x}{10.1111/j.1365-2427.2008.02089.x}

\bibitem[\citeproctext]{ref-Olsson2006}
Olsson K, Stenroth P, Nyström P, Holmqvist N, McIntosh AR, Winterbourn
MJ. 2006. Does natural acidity mediate interactions between introduced
brown trout, native fish, crayfish and other invertebrates in {West
Coast New Zealand} streams? \emph{Biological Conservation} \textbf{130}:
255--267.
https://doi.org/\href{https://doi.org/10.1016/j.biocon.2005.12.019}{10.1016/j.biocon.2005.12.019}

\bibitem[\citeproctext]{ref-Parkyn2001}
Parkyn SM, Collier KJ, Hicks BJ. 2001. {New Zealand} stream crayfish:
Functional omnivores but trophic predators? \emph{Freshwater Biology}
\textbf{46}: 641--652.
https://doi.org/\href{https://doi.org/10.1046/j.1365-2427.2001.00702.x}{10.1046/j.1365-2427.2001.00702.x}

\bibitem[\citeproctext]{ref-Parkyn2009}
Parkyn SM, Meleason MA, Davies‐Colley RJ. 2009. Wood enhances crayfish
({\emph{Paranephrops planifrons}}) habitat in a forested stream.
\emph{New Zealand Journal of Marine and Freshwater Research}
\textbf{43}: 689--700.
https://doi.org/\href{https://doi.org/10.1080/00288330909510034}{10.1080/00288330909510034}

\bibitem[\citeproctext]{ref-Percival2010}
Percival D, Moore P. 2010. Shelter size influences self-assessment of
size in crayfish, {\emph{Orconectes rusticus}}: Consequences for
agonistic fights. \emph{Behaviour} \textbf{147}: 103--119.
https://doi.org/\href{https://doi.org/10.1163/000579509x12512685881053}{10.1163/000579509x12512685881053}

\bibitem[\citeproctext]{ref-Pirtle2012}
Pirtle JL, Eckert GL, Stoner AW. 2012. Habitat structure influences the
survival and predator-prey interactions of early juvenile red king crab
{\emph{Paralithodes camtschaticus}}. \emph{Marine Ecology Progress
Series} \textbf{465}: 169--184.
https://doi.org/\href{https://doi.org/10.3354/meps09883}{10.3354/meps09883}

\bibitem[\citeproctext]{ref-Ramalho2010}
Ramalho RO, McClain WR, Anastacio PM. 2010. An effective and simple
method of temporarily marking crayfish. \emph{Freshwater Crayfish}
\textbf{17}: 57--60.

\bibitem[\citeproctext]{ref-Ranta1993}
Ranta E, Lindström K. 1993. Body size and shelter possession in mature
signal crayfish, {\emph{Pacifastacus leniusculus}}. \emph{Annales
Zoologici Fennici} \textbf{30}: 125--132.

\bibitem[\citeproctext]{ref-Raven2026b}
Raven OV, Burdon FJ, Kusabs IAK, Holmes R, Özkundakci D. 2026b. Littoral
habitat structure influences freshwater crayfish populations in five
volcanic lakes of aotearoa new zealand. \emph{Hydrobiologia}.
https://doi.org/\href{https://doi.org/10.1007/s10750-026-06302-z}{10.1007/s10750-026-06302-z}

\bibitem[\citeproctext]{ref-Raven2026a}
Raven OV, MacNeil C, Lee F, Holmes R, Kusabs IAK, Burdon FJ, Özkundakci
D. 2026a. Fear of the new: A predatory invader provokes stronger
non-consumptive effects on a keystone native prey than a native
predator.
https://doi.org/\href{https://doi.org/10.21203/rs.3.rs-9891519/v1}{10.21203/rs.3.rs-9891519/v1}

\bibitem[\citeproctext]{ref-Reynolds2013}
Reynolds J, Souty-Grosset C, Alastair A. 2013. Ecological roles of
crayfish in freshwater and terrestrial habitats. \emph{Freshwater
Crayfish} \textbf{19}: 197--218.
https://doi.org/\href{https://doi.org/2010.5869/fc.2013.v19-2.197}{2010.5869/fc.2013.v19-2.197}

\bibitem[\citeproctext]{ref-Richman2015}
Richman NI, Böhm M, Adams SB, Alvarez F, Bergey EA, Bunn JJS, Burnham Q,
Cordeiro J, Coughran J, Crandall KA, et al. 2015. Multiple drivers of
decline in the global status of freshwater crayfish (decapoda:
astacidea). \emph{Philosophical Transactions of the Royal Society B:
Biological Sciences} \textbf{370}: 20140060.
https://doi.org/\href{https://doi.org/10.1098/rstb.2014.0060}{10.1098/rstb.2014.0060}

\bibitem[\citeproctext]{ref-Shave1994}
Shave CR, Townsend CR, Crowl TA. 1994. Anti-predator behaviours of a
freshwater crayfish ({\emph{Paranephrops zealandicus}}) to a native and
an introduced predator. \emph{New Zealand Journal of Ecology}
\textbf{18}: 1--10.

\bibitem[\citeproctext]{ref-Spanier1994}
Spanier E. 1994. What are the characteristics of a good artificial reef
for lobsters? \emph{Crustaceana} \textbf{67}: 173--186.

\bibitem[\citeproctext]{ref-Stewart2011}
Stewart C, Tabak J. 2011. Size‐based dominance hierarchies in the {New
Zealand} freshwater crayfish (koura) {\emph{Paranephrops zealandicus}}.
\emph{New Zealand Journal of Marine and Freshwater Research}
\textbf{45}: 281--285.
https://doi.org/\href{https://doi.org/10.1080/00288330.2011.559661}{10.1080/00288330.2011.559661}

\bibitem[\citeproctext]{ref-Streissl2002}
Streissl F, Hödl W. 2002. Habitat and shelter requirements of the stone
crayfish, {\emph{Austropotamobius torrentium}} schrank.
\emph{Hydrobiologia} \textbf{477}: 195--199.
https://doi.org/\href{https://doi.org/10.1023/A:1021094309738}{10.1023/A:1021094309738}

\bibitem[\citeproctext]{ref-Takahashi2019}
Takahashi K, Yamaguchi E, Fujiyama N, Nagayama T. 2019. The effects of
quality of shelters and prior residence on marmorkrebs (marbled
crayfish). \emph{Journal of Experimental Biology}.
https://doi.org/\href{https://doi.org/10.1242/jeb.197301}{10.1242/jeb.197301}

\bibitem[\citeproctext]{ref-Tank1996}
Tank JL, Winterbourn MJ. 1996. Microbial activity and invertebrate
colonisation of wood in a {New Zealand} forest stream. \emph{New Zealand
Journal of Marine and Freshwater Research} \textbf{30}: 271--280.
https://doi.org/\href{https://doi.org/10.1080/00288330.1996.9516714}{10.1080/00288330.1996.9516714}

\bibitem[\citeproctext]{ref-Usio2000}
Usio N, Townsend CR. 2000. Distribution of the {New Zealand} crayfish
{\emph{Paranephrops zealandicus}} in relation to stream
physico‐chemistry, predatory fish, and invertebrate prey. \emph{New
Zealand Journal of Marine and Freshwater Research} \textbf{34}:
557--567.
https://doi.org/\href{https://doi.org/10.1080/00288330.2000.9516957}{10.1080/00288330.2000.9516957}

\bibitem[\citeproctext]{ref-Wehi2017}
Wehi PM, Lord JM. 2017. Importance of including cultural practices in
ecological restoration. \emph{Conservation Biology} \textbf{31}:
1109--1118.
https://doi.org/\href{https://doi.org/10.1111/cobi.12915}{10.1111/cobi.12915}

\bibitem[\citeproctext]{ref-Zhao2024}
Zhao M, Feng G, Wang H, Shen C, Fu Y, Zhang Y, Zhang H, Yao Y, Chen J,
Xu W. 2024. The influence of shelter type and coverage on crayfish
({\emph{Procambarus clarkii}}) predation by catfish ({\emph{Silurus
asotus}}): A controlled environment study. \emph{Animals} \textbf{14}:
1147.
https://doi.org/\href{https://doi.org/10.3390/ani14081147}{10.3390/ani14081147}

\bibitem[\citeproctext]{ref-ZimmerFaust1987}
Zimmer-Faust RK, Spanier E. 1987. Gregariousness and sociality in spiny
lobsters: Implications for den habitation. \emph{Journal of Experimental
Marine Biology and Ecology} \textbf{105}: 57--71.
https://doi.org/\href{https://doi.org/10.1016/s0022-0981(87)80029-1}{10.1016/s0022-0981(87)80029-1}

\end{CSLReferences}




\end{document}
